\PassOptionsToPackage{all}{tcolorbox}
\documentclass[article]{plotweaver}
\usepackage{listings}
\usepackage{lipsum}
\usepackage{graphicx}
\reversemarginpar
\usepackage{geometry}
\geometry{
    paper=letterpaper,
    top=1in,
    bottom=1in,
    left=1in,
    marginparwidth=0.75in,
    marginparsep=0.125in,
    right=1in,
    includeheadfoot,
}

\renewcommand*{\tcbdocupdated}[1]{%
  \textcolor{\maincolor}{\textbf{\textsf{U:}}} #1%
}
\tcbset{
  doc@marginnote/.style={%
    colback=Definition!05!white,
    colframe=Definition,
  },
  doc head={
    interior style={
      fill, 
      color=Gold1!17!white, 
    }
  }
}

\NewDocumentCommand{\dispcmd}{m}{\textcolor{Definition}{\bfseries\texttt{\textbackslash#1}}}
\NewDocumentCommand{\disptxt}{m}{\textbf{\texttt{#1}}}

\title{Plotweaver \LaTeX~Document Class}
\author{}
\date{}

\begin{document}
\maketitle

Documentation for the Plotweaver \LaTeX~document class. See \disptxt{example.tex} for sample usage. More documentation to come.

\section{Icon Commands}
These commands produce commonly used icons. All of the icons are implemented using \disptxt{TikZ}. The icons created by \dispcmd{action}, \dispcmd{doubleAction}, \dispcmd{tripleAction}, \dispcmd{freeAction}, \dispcmd{reaction}, \dispcmd{forever}, and \dispcmd{key} will match the current text color. For example: \textcolor{red}{\action} \textcolor{blue}{\action} \textcolor{green}{\action}. Icons will also respect font size: {\Large\action} \action {\scriptsize\;\action}.

\begin{docCommand}[
    doc name        = action,
    doc description = { \action },
    doc updated     = 2026-08-11,
  ]{}{}
  This command produces the \action icon.
\end{docCommand}

\begin{docCommand}[
    doc name        = doubleAction,
    doc description = { \doubleAction },
    doc updated     = 2026-08-11,
  ]{}{}
  This command produces the \doubleAction icon.
\end{docCommand}

\begin{docCommand}[
    doc name        = tripleAction,
    doc description = { \tripleAction },
    doc updated     = 2026-08-11,
    ]{}{}
  This command produces the \tripleAction icon.
\end{docCommand}

\begin{docCommand}[
    doc name        = freeAction,
    doc description = { \freeAction },
    doc updated     = 2026-08-11,
  ]{}{}
  This command produces the \freeAction icon.
\end{docCommand}

\begin{docCommand}[
    doc name        = reaction,
    doc description = { \reaction },
  ]{}{}
  This command produces the \reaction icon.
\end{docCommand}

\begin{docCommand}[
    doc name        = forever,
    doc description = { \forever },
  ]{}{}
  This command produces the \forever icon.
\end{docCommand}

\begin{docCommand}[
    doc name        = complication,
    doc description = { \complication },
    doc updated     = 2026-08-11,
    doc parameter = \oarg{value},
  ]{}{}
  This command produces the \complication icon. It takes an optional argument for a number, which will be subsequently displayed.
  \begin{dispExample}
  A \complication is complicated. Sometimes complications come with values like \complication[X].
  \end{dispExample}
\end{docCommand}

\begin{docCommand}[
    doc name        = complicationDice,
    doc description = { \complicationDice[x] },
    doc parameter = \oarg{value},
    doc updated     = 2026-08-11,
  ]{}{}
  This command produces the \complicationDice icon. It takes an optional argument for a number, which will be displayed inside of the icon like so: \complicationDice[x]. It is a little larger than \dispcmd{complication}.
  \begin{dispExample}
  A \complicationDice is complicated. Sometimes complications come with values like \complicationDice[x].
  \end{dispExample}
\end{docCommand}

\begin{docCommand}[
    doc name        = opportunity,
    doc description = { \opportunity },
  ]{}{}
  This command produces the \opportunity icon.
\end{docCommand}

\begin{docCommand}[
    doc name        = key,
    doc description = { \key },
    doc updated     = 2026-08-11,
  ]{}{}
  This command produces the \key icon.
\end{docCommand}

\section{Feature Boxes}

\begin{docEnvironment}{featureBox}{\oarg{options}}
  This is a custom tcolorbox and works as such. Under the hood, it is defined by \dispcmd{newtcolorbox}. See the \disptxt{tcolorbox} documentation for more information. 
  \begin{dispExample*}{sidebyside}
    \begin{featureBox}{Feature}
      \lipsum[1][1-3]
  \end{featureBox}
  \end{dispExample*}
\end{docEnvironment}

\begin{docEnvironment}[doc updated=2026-08-11]{statBox}{\marg{heading 1}\marg{content 1}\marg{heading 2}\marg{content 2}}
This is a series of custom boxes that will stay the same height regardless of content length. They are always of equal width. You can use each half of \dispcmd{statBox} separately by using \dispcmd{leftBox} and \dispcmd{rightBox}. These follow the same documentation as \dispcmd{featureBox}.
\begin{dispExample*}{sidebyside}
\statBox{Stat One}{
    10
}%
{Stat Two}{
    20
}
\end{dispExample*}
\end{docEnvironment}

\end{document}